

\chapter*{ABSTRACT}
\phantomsection
\addcontentsline{toc}{chapter}{ABSTRACT}


\bigskip

This project explores a terrain-based localization system designed for the rugged Himalayan landscape of Nepal, where GPS reliability can be affected by signal multipath and ground-level visual data is often limited. The system utilizes Global 30m Elevation Data (Copernicus) (GLO-30) Digital Elevation Model (DEM) data to generate a searchable database of synthetic horizon profiles. For image processing, an efficient segmentation architecture, such as MobileNet-V3-UNet, is proposed to isolate the skyline from user-captured photos. To address potential orientation uncertainty, the project investigates matching strategies—possibly involving Euclidean distance for initial pruning and Dynamic Time Warping (DTW) for refining matches—to account for perspective variations. The system can be evaluated using georeferenced data to assess localization potential and computational feasibility in high-relief environments.

\vspace{1cm}
\noindent \textbf{Keywords}---DEM, Skyline Localization, Himalayan Topography, Computer Vision, MobileNet-V3, DTW.
